\section{Problem}\label{sec:intro:problem}

A large language model (LLM) turns each input token into a sequence of hidden states, one vector per token per layer. Split and hybrid inference systems send those vectors across a trust boundary: a client runs the first blocks, a server runs the rest, and the hidden states travel between them. An honest-but-curious server that holds those states holds the input text as well, because prior work shows that token identity survives in the hidden states and can be decoded back into text~\todo{cite inversion work: vec2text, logit lens, hidden-state inversion}.

Every published inversion attack, however, requires substantial knowledge of the model: its architecture, its tokenizer, and its vocabulary. The attacker trains an inverter on that specific model, or reads its unembedding matrix directly. A server that sees hidden states from a model it cannot identify has no inverter to apply. That gap between \emph{known-model} and \emph{unknown-model} inversion is the subject of this proposal.

A second line of prior work closes the gap in principle. Representation spaces of different LLMs share geometry: Huh et al.~\todo{cite Platonic Representation Hypothesis} argue that models trained on similar data converge toward similar representations, and Jha et al.~\todo{cite vec2vec, Jha 2025} translate embeddings between unknown and known encoders without paired data. If hidden states of an unknown model can be mapped into a space that a known model decodes, then the known model's decoder inverts the unknown model's states.

The proposal builds that space as a \emph{translation hub}: a 512-dimensional space into which the hidden states of three white-box models are mapped, trained so that any model's frozen decoder recovers the input tokens from any model's hub output. An unseen model is then attached to the hub by one affine map, fitted from a small set of calibration pairs. The question is how small that set can be.
